\section{Results and Discussion}
\label{sec:results}

The first question is whether the action set offers anything to route. Table
\ref{tab:complementarity} and Figure~\ref{fig:headroom}
(Appendix~\ref{app:extra}) answer it at the level of individual rows. Gross
complementarity is not scarce. On CodeXGLUE with Laya, \choice{} repairs 14 of
the 60 \noul{} errors, and the per-row oracle reaches 0.783 accuracy against
0.550 for the better fixed action, a headroom of 0.233. The same cell shows why
headroom is a poor promise: \choice{} breaks 15 correct \noul{} answers, so the
net count on the training pool is $-1$ and the net-repair interval spans zero.
The other three binary cells have smaller gross opportunities (headroom
0.067--0.100) and net counts between $-6$ and $+1$. Across the four cells, the
training pool never produces a net repair distinguishable from zero.

Held-out pools do not turn the gross opportunity into a stable net gain. The
CodeXGLUE/Laya cell gains 20.0 points on \texttt{selection\_dev}
($[+1.7,+38.3]$), but the same frozen direction gains only 1.7 points on confirm
($[-18.3,+20.0]$). That gap is a selection-pool effect, not a routing effect:
dev is the pool on which the direction was chosen. FinDVer with Laya is the
exception that tests the narrative. Its confirm net repair is $+13.6$ points
$[+2.0,+24.6]$, and the fixed-action accuracy difference is $+0.136$
$[+0.020,+0.246]$. The corresponding macro-F1 difference, however, is $+0.111$
with an interval that crosses zero, and the better fixed action reverses from
\noul{} on dev to \choice{} on confirm. The cell therefore contains a real
accuracy-level repair signal on confirm, but not a macro-F1 gain, and not a
direction that a train-selected policy could have known to exploit. CodeXGLUE
with Vela moves the other way: the frozen direction is \choice{}-to-\noul{},
the train net is $-6$, and confirm is exactly zero. FinDVer with Vela is negative
on confirm ($-2.3$ points). Complementarity exists in the action set, but it is
sparse, metric-dependent and pool-dependent.

\begin{table*}[t]
\centering
\footnotesize
\caption{Action complementarity with the train-selected direction frozen across
pools ($n=60$ per training cell; FinDVer dev/confirm: 59/44). \emph{Headroom} is the accuracy gain of a per-row
oracle over the better fixed action on \texttt{search\_train};
\emph{repair/harmful} counts rows where the alternative fixes or breaks the
base decision. Net repair is the paired difference in percentage points
(alternative minus base) with a 95\% group-bootstrap interval on dev and
confirm. Headroom is an upper bound, not a deployable gain.}
\label{tab:complementarity}
\begin{tabular}{@{}llrrrr@{}}
\toprule
Task / model & Direction & Headroom & Repair/harmful & Net dev (95\% CI) & Net confirm (95\% CI) \\
\midrule
CodeXGLUE / Laya & Noul$\to$Choice & 0.233 & 14/15 & +20.0 [+1.7, +38.3] & +1.7 [-18.3, +20.0] \\
CodeXGLUE / Vela & Choice$\to$Noul & 0.067 & 4/10 & -3.3 [-15.0, +8.3] & +0.0 [-15.0, +15.0] \\
FinDVer / Laya & Noul$\to$Choice & 0.083 & 5/6 & -6.8 [-23.3, +6.0] & +13.6 [+2.0, +24.6] \\
FinDVer / Vela & Choice$\to$Noul & 0.100 & 6/7 & +0.0 [-9.1, +11.5] & -2.3 [-10.3, +5.0] \\
\bottomrule
\end{tabular}
\end{table*}


The location of the repairs matters as much as their count. Figure
\ref{fig:repair} (Appendix~\ref{app:extra}) plots the repair rate against the
\noul{} probability bin. Repairs concentrate in the middle and high bands,
around $p\in[0.4,0.8]$, and are close to zero in the lowest band. This is the
opposite of what a confidence cascade assumes: a threshold rule switches when
the model is uncertain, which is precisely where \choice{} is least likely to
fix anything, because the rows it can repair are rows where \noul{} was
confident and wrong. It also explains why calibration cannot rescue the router:
recalibration moves probability mass monotonically, but it does not move
repairable rows from the middle of the distribution to the low-confidence tail.

The second question is whether the incremental value of \choice{} is
predictable from the first call. Table~\ref{tab:predictability} reports
out-of-fold ranking quality and the net utility of the frozen dev-selected
policy. Ranking quality varies widely. Logistic regression reaches 0.674 on
CodeXGLUE with Laya and 0.919 on FinDVer with Vela; the tree reaches 0.739 on
CodeXGLUE with Laya. On FinDVer with Laya both predictors are near chance (0.455
and 0.466), unsurprising given that the cell has only five training repairs. The
policy results are weaker than the ranking results. The best dev policy for
CodeXGLUE with Laya is logistic regression with a threshold that switches every
row: dev net $+20.0$ points, confirm net $+1.7$ points and a cost-adjusted
$-3.3$. The FinDVer/Vela cell, despite the highest OOF AUROC, chooses a policy
with dev net $+3.4$ and confirm net $-2.3$ (zero repairs, one regression).
FinDVer with Laya selects a tree that switches one row and gains 2.3 points on
confirm, a point estimate built on a single decision. CodeXGLUE with Vela is
positive on dev ($+3.3$) and negative on confirm ($-1.7$). In no cell does a
high ranking score translate into a confirm net utility that is both positive
and estimated with enough support to matter.

\begin{table*}[t]
\centering
\scriptsize
\caption{Signed incremental-utility prediction. Target
$\Delta_i=\mathbf{1}[\textsc{Choice}_i\text{ correct}]-
\mathbf{1}[\textsc{Noul}_i\text{ correct}]\in\{-1,0,+1\}$. OOF $\rho$ is the
Spearman correlation under group-aware three-fold CV on \texttt{search\_train};
the threshold is selected on \texttt{search\_dev} and frozen for
\texttt{confirm}. Net utility is mean $\Delta$ minus $\lambda=0.05$ per switch,
in points, with a 95\% group-bootstrap interval; the last column is
switches/rows.}
\label{tab:predictability}
\begin{tabular}{@{}lrrrrrr@{}}
\toprule
Task / model & OOF $\rho$ & AUROC repair & AUROC reg. & Dev net & Confirm net (95\% CI) & Switches \\
\midrule
CodeXGLUE / Laya & -0.02 & 0.360 & 0.607 & +21.7 & +0.0 [-18.8, +13.0] & 42/60 \\
CodeXGLUE / Vela & +0.19 & 0.744 & 0.335 & +3.3 & -1.7 [-19.8, +8.5] & 54/60 \\
FinDVer / Laya & +0.13 & 0.458 & 0.707 & +1.7 & +0.0 [-0.3, +0.0] & 1/44 \\
FinDVer / Vela & -0.17 & 0.396 & 0.361 & +3.4 & -2.3 [-7.7, -0.1] & 5/44 \\
\bottomrule
\end{tabular}
\end{table*}


Calibration is the most natural explanation for a failed probability-based
router, so it serves as a control. Table~\ref{tab:calibration} and
Figure~\ref{fig:calibration} in Appendix~\ref{app:extra} show that raw
probabilities are poorly calibrated on confirm: ECE ranges from 0.090 to 0.290
and Brier scores from 0.218 to 0.287. Temperature scaling, fitted on
\texttt{search\_train}, reduces ECE in three of four cells (to 0.045--0.169) and
Brier in all four. Routing is unaffected: raw probabilities give confirm net
repairs of $-5.0$, $0.0$, $+2.3$ and $-2.3$ points, temperature-scaled ones give
$-1.7$, $+1.7$, $0.0$ and $0.0$, as the invariance argument in
Section~\ref{sec:method} predicts. The repair AUROC makes the point directly:
0.791, 0.444, 0.336 and 0.095 for the four cells; in the last two the
probability ranks repairable rows \emph{below} chance. A calibrated probability
of the current answer is not a calibrated probability of the alternative
action's value.

The third question is whether search converts the available signal into a
deployable policy. Table~\ref{tab:search} summarises the comparison. GA versus
random search is the cleanest test of optimiser value: with five seeds, four
task--model cells and three pools, there are 60 paired macro-F1 comparisons, and
none has an interval that excludes zero. GA versus fixed switch-only has five
significant positive cells, all on the FinDVer/Vela training pool and all with
an effect near $+0.12$; on dev and confirm, none of the 40 comparisons is
significant. The pre-registered deployment comparison---train-selected GA
against fixed switch-only on confirm---is null in every seed-level cell. ES
versus random search is null in all 30 comparisons; ES versus fixed switching
has three significant training cells and one significant confirm cell
(CodeXGLUE/Laya, seed 17, $+0.176$), the only confirm-level positive result, and
it is not replicated across seeds. NSGA-II deployment is significantly worse
than the fixed baseline in five dev cells and one confirm cell, with no positive
cell. The search problem has no stable held-out signal for any optimiser to
find.

\begin{table}[!b]
\centering
\small
\caption{Search comparisons under the pre-specified protocol. Counts are
paired group-bootstrap 95\% intervals that exclude zero; ``positive'' means the
evolved policy beats the stated baseline. GA and ES rows use macro-F1 on all
three pools with five seeds; the NSGA-II row uses accuracy on 12 dev and 12
confirm cells.}
\label{tab:search}
\begin{tabular}{@{}lrrl@{}}
\toprule
Comparison & Pos. & Neg. & Reading \\
\midrule
GA vs.\ random & 0 & 0 & no advantage (60) \\
GA vs.\ fixed & 5 & 0 & train-only (60) \\
ES vs.\ random & 0 & 0 & no advantage (30) \\
ES vs.\ fixed & 4 & 0 & one confirm (30) \\
NSGA-II vs.\ fixed & 0 & 6 & 5 dev, 1 confirm \\
\bottomrule
\end{tabular}
\end{table}


The trajectory diagnostics explain why more search does not help. In 15 of 21
logged GA/ES runs, best-so-far fitness does not improve after generation zero.
The neutral-mutation rate is 0.68--0.73, so most mutations produce a candidate
with exactly the same training score. The train--dev rank correlation is $-0.56$
and $-0.51$ for the two CodeXGLUE cells, meaning that candidates which look
better on train tend to look worse on dev, and $+0.44$ and $+0.99$ for the
FinDVer cells. Raising the budget from 24 to 48 on CodeXGLUE with Vela does not
close the gap: the GA evaluates only 22, 17 and 23 unique genomes under a cap of
48, while random search uses all 48. The search is not merely under-tuned; it is
under-informed.

The boundary checks ask whether these conclusions are artefacts of one pool or
task. Appendix~\ref{app:boundary} reports the transfer table and fault map.
Transfer is uniformly unhelpful: of 32 source--target by split units, six have a
significant paired difference against the target's fixed switch policy, and all
six are negative (two diagonal, four off-diagonal). The fault map is more
nuanced: of 15 eligible cells, 13 are confirmed on dev and 10 on confirm. The
highest confirm severities are on ProcessBench (Vela/mid 0.833, Laya/mid 0.875,
Vela/low 0.936, Laya/high 0.895). The mid-confidence cells are not simply the
low-confidence tail: a model can be moderately confident and still fail, and the
second representation fails on the same rows, so no gate on the first action's
probability can create a repair.

Taken together, the three questions describe a routing problem that fails before
the optimiser is asked to solve it: RQ1 finds gross but no stable net
complementarity, RQ2 finds ranking signal without reliable policy value, and RQ3
finds no held-out advantage because the landscape is flat and the training
ranking does not predict the dev ranking where it matters. Calibration improves
a different quantity: a model can report better probabilities for the question
it answered and still have nothing to say about whether a different question
would be answered better.

The practical implication is a change in the order of work. Before building a
router, measure whether the action set has net complementarity on held-out
groups and whether the candidate features rank repairable rows above chance. If
either check fails, a stronger optimiser will not help; if both pass, the
optimiser comparison can be run under a matched budget with a pre-registered
confirm comparison.
